%!TEX root = ../../../main.tex
%
% §7.3 07c_fccd_advection.tex の FCCD 行列 M^FCCD インライン再定義
%   退避: CHK-193 Phase D-2（2026-04-24）
%
% 退避理由:
%   §4.5 04e_fccd.tex L118 の boxed 一次定義（eq:fccd_composite_matrix）と
%   同式を §7c 本文中でインライン再掲していた．
%   §4e 側に cross-ref が既にあるため，本体の数式展開は冗長．
%
% 本編の代替:
%   paper/sections/07c_fccd_advection.tex L13 付近および L152 付近を
%   §\ref{sec:fccd_def} / 式~\eqref{eq:fccd_composite_matrix} への参照のみに圧縮．
%
% 元の箇所と本文（参考）:

% L13 付近（subsection 冒頭の説明文）:
SP-D \cite{SPD} は \S\ref{sec:fccd_def} で導入した面中心 FCCD 演算子
$\mathbf{M}^{\FCCD} = \mathbf{D}_1 - \mathbf{D}_2 \mathbf{S}_{\CCD}$ を
速度移流項へ拡張し，

% L152 付近（BF 整合性定理の演算子説明）:
\begin{theorem}[FCCD Option B の BF 整合性]\label{thm:fccd_adv_bf}
$\nabla p$（FCCD），$\sigma\kappa\nabla\psi$（FCCD 面形），
$\nabla\cdot(\mathbf{u}\otimes\mathbf{u})$（FCCD Option B）の
3 項すべてを\textbf{共通の面中心演算子}
$\mathbf{M}^{\FCCD} = \mathbf{D}_1 - \mathbf{D}_2\mathbf{S}_{\CCD}$ で
離散化すれば，静止状態 $\mathbf{u}\equiv 0$ での離散運動量バランスの
BF 残差は
...
